% Options for packages loaded elsewhere
\PassOptionsToPackage{unicode}{hyperref}
\PassOptionsToPackage{hyphens}{url}
\documentclass[
  a4paper,
]{article}
\usepackage{xcolor}
\usepackage[margin=25mm]{geometry}
\usepackage{amsmath,amssymb}
\setcounter{secnumdepth}{-\maxdimen} % remove section numbering
\usepackage{iftex}
\ifPDFTeX
  \usepackage[T1]{fontenc}
  \usepackage[utf8]{inputenc}
  \usepackage{textcomp} % provide euro and other symbols
\else % if luatex or xetex
  \usepackage{unicode-math} % this also loads fontspec
  \defaultfontfeatures{Scale=MatchLowercase}
  \defaultfontfeatures[\rmfamily]{Ligatures=TeX,Scale=1}
\fi
\usepackage{lmodern}
\ifPDFTeX\else
  % xetex/luatex font selection
\fi
% Use upquote if available, for straight quotes in verbatim environments
\IfFileExists{upquote.sty}{\usepackage{upquote}}{}
\IfFileExists{microtype.sty}{% use microtype if available
  \usepackage[]{microtype}
  \UseMicrotypeSet[protrusion]{basicmath} % disable protrusion for tt fonts
}{}
\makeatletter
\@ifundefined{KOMAClassName}{% if non-KOMA class
  \IfFileExists{parskip.sty}{%
    \usepackage{parskip}
  }{% else
    \setlength{\parindent}{0pt}
    \setlength{\parskip}{6pt plus 2pt minus 1pt}}
}{% if KOMA class
  \KOMAoptions{parskip=half}}
\makeatother
\setlength{\emergencystretch}{3em} % prevent overfull lines
\providecommand{\tightlist}{%
  \setlength{\itemsep}{0pt}\setlength{\parskip}{0pt}}
\usepackage{bookmark}
\IfFileExists{xurl.sty}{\usepackage{xurl}}{} % add URL line breaks if available
\urlstyle{same}
\hypersetup{
  pdftitle={Emergent Condensate{,} Center of Mass{,} and Mass-Like Measure from an Anonymous Complex Standing-Wave Vacuum: An Anonymous Fourier--Mellin Readout of Condensation Strength{,} Finite Extent{,} and Vacuum-State Gauge Distortion},
  pdfauthor={Noriaki Kihara (Independent Researcher), ORCID 0009-0004-6753-4020},
  hidelinks,
  pdfcreator={LaTeX via pandoc}}

\title{Emergent Condensate, Center of Mass, and Mass-Like Measure from
an Anonymous Complex Standing-Wave Vacuum: An Anonymous Fourier--Mellin
Readout of Condensation Strength, Finite Extent, and Vacuum-State Gauge
Distortion}
\author{Noriaki Kihara (Independent Researcher), ORCID
0009-0004-6753-4020}
\date{September 3, 2026 --- v3.0 --- Version DOI
10.5281/zenodo.22282690}

\begin{document}
\maketitle

\subsection{Abstract}\label{abstract}

Without assuming localized particles, mass, a center of mass, or
background spacetime from the start, this paper begins with an anonymous
set of complex standing waves and the complex square-zero closure

\[
\sum_n z_n^2=0,
\]

and constructs a method to read out, when a localized condensate forms
inside the homogeneous vacuum state field, the condensate's relative
center of mass, finite extent, condensation strength, and the distortion
of the surrounding vacuum-state gauge, without introducing external
coordinates or particle labels.

In the previous paper {[}K1{]}, relative scales, phase periods,
logarithmic spatial and temporal readouts, and a post-mapping
\[\lambda\nu\]-type invariant were constructed from the anonymous
complex zero-closure system. In this paper, that homogeneous vacuum
state field is taken as a ``sea of many complex standing waves'', and we
treat the case where a localized condensation distinguishable from the
homogeneous background arises within it.

For a complex number

\[
z_n=r_n e^{i\theta_n},
\]

using the logarithmic amplitude

\[
X_n=\ln r_n
\]

and the phase \[\theta_n\], we define a Fourier--Mellin-type transform
of the empirical measure of the anonymous set,

\[
\mathcal A(s,m)
=
\frac{1}{N}
\sum_{n=1}^{N}
\exp\left[
-i\left(
s\ln r_n+m\theta_n
\right)
\right].
\]

This transform is invariant under rearrangements of the waves, and the
common amplitude scale and common phase rotation appear only as overall
phase factors of the transform, so that its intensity

\[
P(s,m)=|\mathcal A(s,m)|^2
\]

is invariant under anonymity, the common scale, and the common phase.

Writing the normalized transform of a localized condensate as
\[\chi_C\], its logarithm expands near the transform origin as

\[
\ln\chi_C(\mathbf q)
=
-i\mathbf q\cdot\boldsymbol\mu_C
-\frac12
\mathbf q^{\mathsf T}\Sigma_C\mathbf q
+O(|\mathbf q|^3),
\]

where \[\boldsymbol\mu_C\] is the condensate's center of mass and
\[\Sigma_C\] is the covariance matrix expressing its finite extent and
shape. Therefore the center of mass can be read from the first
derivative of the spectral phase, and the extent from the second
derivative of the spectral amplitude. The absolute center of mass,
however, is erased by the common gauge transformation; what can be
physically read is the relative center of mass with respect to another
condensate or the homogeneous vacuum background.

Furthermore, introducing the difference

\[
\delta\mathcal A
=
\mathcal A-\mathcal A_0
\]

between the pre-condensation homogeneous vacuum spectrum
\[\mathcal A_0\] and the post-condensation spectrum \[\mathcal A\], the
existence of a condensate can be defined as a local inhomogeneity of the
vacuum state field. Since the inverse-transformed \[\delta\rho\] can
extend outside the condensation center, the condition under which
localized condensation deforms the surrounding vacuum-state gauge can be
verified directly. Without placing mass as an external parameter, this
paper proposes a ``mass-like measure'' that obtains the condensation
strength, a stable relative center of mass, a finite second moment, and
the background gauge deformation from one and the same change of state.

In Appendix A, we further generalize the previous paper's light-like
special solution to \[P=P_0L^\alpha\] and show that the necessary and
sufficient condition for \[\lambda\nu\] to be constant across all shell
pairs is this power relation. Matching the readout units, \[\alpha<1\],
\[\alpha=1\], and \[\alpha>1\] correspond to spacelike, null, and
timelike, respectively. We also construct a rational-exponent shell
family that simultaneously satisfies integer cycles, exact per-shell
square-zero closure, and overall absolute convergence. The connection
claiming that localized condensation actually generates
\[\alpha_{\mathrm{eff}}>1\] is separated as an underived hypothesis.

The claim of this paper is not that real particle masses or the Einstein
equations have already been derived. What is shown is an explicit
readout framework in which, for a square-zero-closed anonymous complex
standing-wave vacuum, a localized condensate can be detected anonymously
and its relative center of mass, finite extent, condensation strength,
and the distortion of the background state field can all be defined from
the same harmonic analysis.

\textbf{Keywords}: anonymous complex numbers, square-zero closure,
standing-wave vacuum, localized condensate, center of mass, mass,
Fourier--Mellin transform, characteristic function, gauge-invariant
readout, emergent matter

\begin{center}\rule{0.5\linewidth}{0.5pt}\end{center}

\subsection{1. Problem Setting --- How to Read an ``Object'' out of a
Homogeneous
Vacuum}\label{problem-setting-how-to-read-an-object-out-of-a-homogeneous-vacuum}

In the previous paper {[}K1{]}, without first placing space, time,
wavelength, or frequency, we investigated what can be read out
relatively from an anonymous set of complex numbers

\[
\mathcal Z
=
\{z_1,z_2,\ldots\}/S
\]

and the square-zero closure

\[
Q(\mathcal Z)
=
\sum_n z_n^2
=
0.
\tag{1}
\]

Equation (1) is preserved under any common complex scale

\[
z_n\mapsto \gamma z_n,
\qquad
\gamma\in\mathbb C^\times,
\]

since

\[
Q\mapsto \gamma^2Q=0.
\]

Therefore the closure condition itself selects neither an absolute
amplitude scale nor an absolute phase.

In this paper, such a state is read as a homogeneous vacuum state field
consisting of many complex standing waves. What matters is that this is
not a ``vacuum field existing inside spacetime'' but a state set prior
to the readout spacetime.

In a homogeneous vacuum, no special center is placed anywhere. Therefore
the vacuum itself has no center of mass intrinsic to an object.

The question of this paper is the following.

\begin{quote}
When a part of the anonymous sea of complex standing waves condenses and
forms a localized structure distinguishable from the background, can its
center of mass, extent, strength, and the distortion of the surrounding
background field be read out without giving external coordinates,
particle numbers, or known masses?
\end{quote}

Organizing the necessary conditions first, a candidate material
condensate must satisfy at least the following.

\begin{enumerate}
\def\labelenumi{\arabic{enumi}.}
\tightlist
\item
  It does not break the square-zero closure of the whole system.
\item
  It is statistically and spectrally distinguishable from the
  homogeneous background.
\item
  It has a stable relative center of mass.
\item
  It has a finite extent that is neither zero nor infinite.
\item
  It can induce a deformation of the background state outside the
  condensation center.
\item
  At large distance, or under sufficiently coarse readout, it returns to
  the original homogeneous vacuum.
\end{enumerate}

A structure satisfying these six conditions is called a ``localized
condensate'' in this paper.

\begin{center}\rule{0.5\linewidth}{0.5pt}\end{center}

\subsection{2. Scale Invariance of the Square-Zero Closure and
Interference
Closure}\label{scale-invariance-of-the-square-zero-closure-and-interference-closure}

Writing each complex number as

\[
z_n=a_n+i b_n,
\]

we have

\[
z_n^2
=
a_n^2-b_n^2
+
2i a_n b_n.
\]

Therefore equation (1) splits into

\[
\sum_n(a_n^2-b_n^2)=0
\tag{2}
\]

and

\[
\sum_n a_n b_n=0.
\tag{3}
\]

In particular, equation (3) requires that the cross terms of the complex
components, that is, the total of the interference terms appearing in
the square readout, close.

Transforming all components by a common real scale \[\kappa>0\],

\[
a_n\mapsto \kappa a_n,
\qquad
b_n\mapsto \kappa b_n,
\]

we get

\[
\sum_n(\kappa a_n)(\kappa b_n)
=
\kappa^2
\sum_n a_n b_n
=
0.
\]

Therefore the zero closure fixes no scale, and has the property

\[
\boxed{
\text{self-consistency does not depend on any absolute scale}
}
\]

This paper places importance on this property. Mass and length are not
inserted into the foundational equation as bare absolute parameters;
they are treated as relative readouts arising after condensation.

\begin{center}\rule{0.5\linewidth}{0.5pt}\end{center}

\subsection{3. The Homogeneous Vacuum State Gauge
Field}\label{the-homogeneous-vacuum-state-gauge-field}

\subsubsection{3.1 Definition of the
vacuum}\label{definition-of-the-vacuum}

Write the vacuum state as

\[
\mathcal V_0
=
\{z_n^{(0)}\}
\]

satisfying

\[
\sum_n
\left(z_n^{(0)}\right)^2
=
0.
\tag{4}
\]

Furthermore, the vacuum is assumed to have no localized special point
with respect to the anonymous readout adopted. In this sense it is
called ``homogeneous''.

The state gauge here means that no absolute common amplitude or common
phase is physically fixed. That is, states related by

\[
z_n^{(0)}
\mapsto
\gamma z_n^{(0)}
\]

within the same closure class belong to the same gauge orbit from the
viewpoint of relative readout.

\subsubsection{3.2 Vacuum and matter are not separate foundational
variables}\label{vacuum-and-matter-are-not-separate-foundational-variables}

In this paper, the vacuum and the condensate are not introduced as
different kinds of entities.

Write the post-condensation state as

\[
\mathcal V
=
\{z_n\},
\]

maintaining

\[
\sum_n z_n^2=0.
\tag{5}
\]

Material condensation refers to the case where, among the
self-consistent rearrangements from \[\mathcal V_0\] to \[\mathcal V\],
a localized correlation structure arises.

We therefore adopt the one-layer structure

\[
\boxed{
\text{vacuum}
\longrightarrow
\text{localized rearrangement of the same state degrees of freedom}
\longrightarrow
\text{condensate}
}
\]

\begin{center}\rule{0.5\linewidth}{0.5pt}\end{center}

\subsection{4. The Empirical Measure of the Anonymous
Set}\label{the-empirical-measure-of-the-anonymous-set}

Since wave numbers carry no physical meaning, readout operations must
not depend on rearrangements of the indices.

For a finite set, define the empirical measure

\[
\mu
=
\frac1N
\sum_{n=1}^{N}
\delta_{z_n}.
\tag{6}
\]

Equation (6) is invariant under rearrangements of the set.

Represent the complex numbers in polar form

\[
z_n
=
r_n e^{i\theta_n},
\qquad
r_n>0.
\]

Zeros have no logarithmic amplitude, so if needed \[z=0\] is treated
separately as a limit point, and the transform is defined on
\[\mathbb C^\times\].

The multiplicative group of nonzero complex numbers decomposes as

\[
\mathbb C^\times
\simeq
\mathbb R_+
\times
S^1.
\]

Therefore it is natural to combine harmonic analysis in the amplitude
direction \[r\in\mathbb R_+\] with harmonic analysis in the phase
direction \[\theta\in S^1\].

Adopting the logarithmic scale readout of the previous paper {[}K1{]},

\[
X=\ln r,
\]

multiplicative amplitude ratios are carried to additive differences.

\begin{center}\rule{0.5\linewidth}{0.5pt}\end{center}

\subsection{5. The Anonymous Fourier--Mellin
Transform}\label{the-anonymous-fouriermellin-transform}

\subsubsection{5.1 Definition}\label{definition}

For the anonymous complex set, define

\[
\boxed{
\mathcal A(s,m)
=
\frac1N
\sum_{n=1}^{N}
\exp
\left[
-i
\left(
s\ln r_n
+
m\theta_n
\right)
\right]
}
\tag{7}
\]

where

\[
s\in\mathbb R,
\qquad
m\in\mathbb Z.
\]

The phase direction is the circle \[S^1\], giving a Fourier series; the
amplitude direction is the multiplicative group \[\mathbb R_+\], giving
a Mellin-type transform.

Using the empirical measure, equation (7) can be written as

\[
\mathcal A(s,m)
=
\int_{\mathbb C^\times}
|z|^{-is}
\left(
\frac{z}{|z|}
\right)^{-m}
d\mu(z).
\tag{8}
\]

\subsubsection{5.2 Permutation invariance}\label{permutation-invariance}

For any permutation \[\pi\], replacing

\[
z_n\mapsto z_{\pi(n)}
\]

does not change the sum. Therefore \[\mathcal A\] preserves anonymity.

\subsubsection{5.3 Common scale
transformation}\label{common-scale-transformation}

If

\[
r_n
\mapsto
\kappa r_n,
\]

then

\[
\ln r_n
\mapsto
\ln r_n+\ln\kappa.
\]

Hence

\[
\mathcal A(s,m)
\mapsto
 e^{-is\ln\kappa}
\mathcal A(s,m).
\tag{9}
\]

Therefore the intensity

\[
\boxed{
P(s,m)
=
|\mathcal A(s,m)|^2
}
\tag{10}
\]

is invariant under the common scale.

\subsubsection{5.4 Common phase rotation}\label{common-phase-rotation}

If

\[
\theta_n
\mapsto
\theta_n+\phi,
\]

then

\[
\mathcal A(s,m)
\mapsto
 e^{-im\phi}
\mathcal A(s,m).
\tag{11}
\]

Therefore equation (10) is invariant under the common phase as well.

From the above, \[P(s,m)\] is an anonymous spectrum that does not depend
on wave numbering, the absolute amplitude scale, or the absolute phase.

That Fourier--Mellin-type analysis can be used for representations
invariant under scale transformations and rotations is itself known
{[}E1{]}. The original point of this paper is to apply it not to images
or existing spatial coordinates, but to a pre-spacetime anonymous
complex state set, and to connect it to the material readout of
condensates.

\begin{center}\rule{0.5\linewidth}{0.5pt}\end{center}

\subsection{6. Definition of a
Condensate}\label{definition-of-a-condensate}

Let \[\mu_0\] be the vacuum state and \[\mu\] the post-condensation
measure.

Set the formal difference

\[
\delta\mu
=
\mu-\mu_0.
\tag{12}
\]

A condensate is a state in which \[\delta\mu\] contains a correlation
component that is separable from the background and has finite width.

Here we write the condensate component as \[\mu_C\] and its total weight
as

\[
M_0
=
\int d\mu_C.
\tag{13}
\]

The quantity \[M_0\] is not yet called a physical mass; it is called the
``condensation participation weight'' or the ``zeroth condensation
moment''.

Define the normalized condensate measure

\[
d\nu_C
=
\frac{d\mu_C}{M_0}.
\tag{14}
\]

\begin{center}\rule{0.5\linewidth}{0.5pt}\end{center}

\subsection{7. The Center of Mass Appears as the First Derivative of the
Spectral
Phase}\label{the-center-of-mass-appears-as-the-first-derivative-of-the-spectral-phase}

\subsubsection{7.1 Derivation on the one-dimensional logarithmic
scale}\label{derivation-on-the-one-dimensional-logarithmic-scale}

First fix the phase direction and look only at the logarithmic scale

\[
X=\ln r.
\]

Let the normalized characteristic function of the condensate be

\[
\chi_C(s)
=
\int
 e^{-isX}
d\nu_C(X).
\tag{15}
\]

At the origin,

\[
\chi_C(0)=1.
\]

Differentiating,

\[
\left.
\frac{d\chi_C}{ds}
\right|_{s=0}
=
-i
\int X\,d\nu_C
=
-iX_C,
\tag{16}
\]

where

\[
\boxed{
X_C
=
\int X\,d\nu_C
}
\tag{17}
\]

is defined as the condensate's center of mass.

Using the logarithm of the characteristic function,

\[
\boxed{
X_C
=
i
\left.
\frac{d}{ds}
\ln\chi_C(s)
\right|_{s=0}
}
\tag{18}
\]

Written with the phase sign convention,

\[
\boxed{
X_C
=
-
\left.
\frac{d}{ds}
\arg\chi_C(s)
\right|_{s=0}
}
\tag{19}
\]

This shows that the central position of the condensate can be read as
the first-order slope of the spectral phase.

\subsubsection{7.2 The absolute center of mass is not read
out}\label{the-absolute-center-of-mass-is-not-read-out}

Shifting the whole condensate by

\[
X\mapsto X+X_0,
\]

we get

\[
\chi_C(s)
\mapsto
 e^{-isX_0}\chi_C(s).
\tag{20}
\]

Therefore \[|\chi_C|\] does not change, but the phase slope changes by
\[X_0\].

This is not a defect. Since the underlying system has no absolute scale,
the absolute center of mass should not be physically fixed either.

For two condensates \[A,B\], taking

\[
\frac{\chi_A(s)}{\chi_B(s)},
\]

we obtain

\[
\boxed{
X_A-X_B
=
-
\left.
\frac{d}{ds}
\arg
\frac{\chi_A(s)}{\chi_B(s)}
\right|_{s=0}
}
\tag{21}
\]

Therefore what can be read out is not the absolute center of mass but
the relative center of mass.

\begin{center}\rule{0.5\linewidth}{0.5pt}\end{center}

\subsection{8. Finite Extent Appears as the Second Derivative of the
Spectral
Amplitude}\label{finite-extent-appears-as-the-second-derivative-of-the-spectral-amplitude}

Let the condensate's variance be

\[
\sigma_X^2
=
\int
(X-X_C)^2
d\nu_C.
\tag{22}
\]

By the cumulant expansion of the characteristic function {[}E2{]},

\[
\ln\chi_C(s)
=
-isX_C
-
\frac12
s^2\sigma_X^2
+
O(s^3).
\tag{23}
\]

Therefore

\[
\boxed{
\sigma_X^2
=
-
\left.
\frac{d^2}{ds^2}
\Re\ln\chi_C(s)
\right|_{s=0}
}
\tag{24}
\]

This gives the clear hierarchy

\[
\boxed{
\text{first-order spectral phase}
\longrightarrow
\text{center of mass}
}
\]

\[
\boxed{
\text{second-order spectral decay}
\longrightarrow
\text{extent}
}
\]

Requiring

\[
0<\sigma_X^2<\infty
\tag{25}
\]

as the minimal condition of a material condensate, we can define a
finite-size object that is neither a perfect point singularity nor an
infinitely diffused background.

\begin{center}\rule{0.5\linewidth}{0.5pt}\end{center}

\subsection{9. Two-Variable Expansion Including the Phase/Time
Direction}\label{two-variable-expansion-including-the-phasetime-direction}

Including the phase direction, set

\[
Y=
\begin{pmatrix}
X\\
T
\end{pmatrix}.
\]

Here \[T\] is the time readout obtained from the phase period or
relative winding number in the previous paper {[}K1{]}.

With the dual variable

\[
q=
\begin{pmatrix}
s\\
\omega
\end{pmatrix},
\]

we have

\[
\chi_C(q)
=
\int
 e^{-iq^{\mathsf T}Y}
d\nu_C(Y).
\tag{26}
\]

Near the origin,

\[
\boxed{
\ln\chi_C(q)
=
-iq^{\mathsf T}\mu_C
-
\frac12
q^{\mathsf T}
\Sigma_C
q
+
O(|q|^3)
}
\tag{27}
\]

where

\[
\mu_C
=
\begin{pmatrix}
X_C\\
T_C
\end{pmatrix}
\]

is the center of mass, and

\[
\Sigma_C
=
\begin{pmatrix}
\sigma_X^2 & C_{XT}\\
C_{XT} & \sigma_T^2
\end{pmatrix}
\]

is the covariance matrix.

Therefore the cross term

\[
\boxed{
C_{XT}
=
-
\left.
\frac{\partial^2}
{\partial s\,\partial\omega}
\Re\ln\chi_C
\right|_{q=0}
}
\tag{28}
\]

can also be read out.

With this, not only the mere radius of the condensate but also the
coupling between the spatial scale and the temporal phase can be
measured.

\begin{center}\rule{0.5\linewidth}{0.5pt}\end{center}

\subsection{10. Generalization to a 3-Space + Time
Readout}\label{generalization-to-a-3-space-time-readout}

\textbf{Assumption 10.1 (existence of a four-dimensional relative
readout)}\\
We assume that the spacetime readout of the previous paper {[}K1{]} and
the relative geometry of multiple complex planes can be extended so that
four-dimensional relative readout coordinates can be constructed for the
condensate. This is not a proposition newly proved in this paper. The
readout explicitly constructed in the previous paper is first of all
\[1+1\] dimensional, and the four-dimensionalization below is a
conditional construction valid when that extension holds.

Concretely, suppose readout coordinates

\[
Y_n
=
\begin{pmatrix}
X_n\\
Y_n\\
Z_n\\
T_n
\end{pmatrix}
\]

have been constructed for the condensate.

What matters here is that these are not given from the start as a
background lattice; they are relative readouts obtained from complex
state relations.

Consequently, the four-dimensional moment readout of equations
(29)--(31), and the four-dimensional mass-like interpretation of Section
11 that uses it, are conditional on Assumption 10.1. On the other hand,
the one-dimensional and two-variable moment identities derived in
Sections 7--9 themselves do not depend on this assumption.

With the four-dimensional dual variable

\[
q
=
\begin{pmatrix}
k_x\\
k_y\\
k_z\\
\omega
\end{pmatrix},
\]

set

\[
\boxed{
\chi_C(q)
=
\frac{1}{M_0}
\int
\exp
\left[
-i
\left(
k_xX+k_yY+k_zZ-\omega T
\right)
\right]
d\mu_C
}
\tag{29}
\]

The sign of the time term is chosen opposite for consistency with the
Lorentz-type readout, but the moment extraction itself does not depend
on the sign convention.

The logarithmic expansion of equation (29) is

\[
\boxed{
\ln\chi_C(q)
=
-iq^{\mathsf T}\mu_C
-
\frac12q^{\mathsf T}\Sigma_C q
+
O(|q|^3)
}
\tag{30}
\]

where

\[
\boxed{
\mu_C
=
(X_C,Y_C,Z_C,T_C)
}
\tag{31}
\]

is the four-dimensional relative center of mass and \[\Sigma_C\] is the
\[4\times4\] covariance matrix.

From the eigenvalues and eigenvectors of the spatial \[3\times3\] part,
the principal radii, ellipticity, principal axes, and anisotropy can be
read out.

A single spectral operation therefore yields the condensation weight,
center of mass, extent, shape, and space-time correlation in a unified
way.

\begin{center}\rule{0.5\linewidth}{0.5pt}\end{center}

\subsection{11. The Mass-Like Measure}\label{the-mass-like-measure}

\subsubsection{11.1 Mass is not placed as a foundational
parameter}\label{mass-is-not-placed-as-a-foundational-parameter}

In this paper the physical mass \[m\] is not introduced as an initial
condition.

The homogeneous vacuum has no localized center.

Only when condensation occurs do the condensation participation weight,
the relative center of mass, the finite extent, and the background
distortion become definable simultaneously.

We therefore define the minimal data of a mass-like object as

\[
\boxed{
\mathfrak M_C
=
\left(
M_0,
\mu_C,
\Sigma_C
\right)
}
\tag{32}
\]

This is not yet an inertial mass in kilograms.

The ``mass-like measure'' of this paper is a state quantity expressing
the condensation strength that creates a localized center and finite
width relative to the homogeneous vacuum.

\subsubsection{11.2 Condensation weight as the zeroth
moment}\label{condensation-weight-as-the-zeroth-moment}

With the unnormalized condensation spectrum

\[
\Delta\mathcal A_C(q)
=
\int
 e^{-iq^{\mathsf T}Y}
d\mu_C(Y),
\tag{33}
\]

we have

\[
\boxed{
\Delta\mathcal A_C(0)
=
M_0
}
\tag{34}
\]

Therefore the hierarchy

\[
\boxed{
\begin{aligned}
0\text{th moment}
&\longrightarrow
\text{condensation weight},\\
1\text{st moment}
&\longrightarrow
\text{center of mass},\\
2\text{nd moment}
&\longrightarrow
\text{extent and shape}
\end{aligned}
}
\tag{35}
\]

holds.

This structure gives a minimal candidate for measuring mass from the
state distribution of the condensate itself, rather than assigning it as
an external scalar.

\begin{center}\rule{0.5\linewidth}{0.5pt}\end{center}

\subsection{12. Cross Phase against the Homogeneous Vacuum and Relative
Position}\label{cross-phase-against-the-homogeneous-vacuum-and-relative-position}

The spectral intensity

\[
|\mathcal A|^2
\]

alone loses the absolute position information.

Therefore, reading the condensate's center of mass requires a cross
phase with a reference.

Let \[\mathcal A_0(q)\] be the pre-condensation vacuum and
\[\mathcal A(q)\] the post-condensation spectrum.

Consider the difference spectrum

\[
\boxed{
\delta\mathcal A(q)
=
\mathcal A(q)-\mathcal A_0(q)
}
\tag{36}
\]

Alternatively, in regions where the reference spectrum is nonzero, the
cross spectrum

\[
\boxed{
C(q)
=
\mathcal A(q)
\mathcal A_0(q)^*
}
\tag{37}
\]

may be used.

The phase obtained here is not a position relative to an external origin
but a relative phase shift with respect to the homogeneous vacuum gauge.

This structure is consistent with the basic principle of this paper:
read only relations, not absolute quantities.

\begin{center}\rule{0.5\linewidth}{0.5pt}\end{center}

\subsection{13. Distortion of the Surrounding Vacuum State Gauge by a
Condensate}\label{distortion-of-the-surrounding-vacuum-state-gauge-by-a-condensate}

\subsubsection{13.1 Localized condensation and global
closure}\label{localized-condensation-and-global-closure}

The whole system always maintains

\[
\sum_n z_n^2=0.
\]

Formally dividing the condensate region \[C\] from its exterior \[E\],

\[
Q_C
=
\sum_{n\in C}z_n^2,
\]

\[
Q_E
=
\sum_{n\in E}z_n^2,
\]

we have

\[
\boxed{
Q_C+Q_E=0
}
\tag{38}
\]

In general there is no need for the condensate part alone to satisfy
\[Q_C=0\].

If the localized condensation has

\[
Q_C\neq0,
\]

the global closure necessarily requires

\[
\boxed{
Q_E=-Q_C
}
\tag{39}
\]

Therefore, if the closure contribution becomes biased by local
condensation, the surrounding state set must rearrange compensatingly.

This is a candidate for producing

\[
\boxed{
\text{localized condensation}
\longrightarrow
\text{compensating deformation of the background state}
}
\]

without additional axioms.

\subsubsection{13.2 Reducible closure and the limits of the
residual-compensation
mechanism}\label{reducible-closure-and-the-limits-of-the-residual-compensation-mechanism}

The explicit shell special solution of the previous paper {[}K1{]} has a
reducible structure in which each shell individually satisfies the
square-zero closure. For this reason, when the condensate region \[C\]
is composed only of complete closure shells,

\[
Q_C=0
\]

holds, and the \textbf{mechanism by which the exterior compensates a
nonzero closure residual} of equation (39) does not operate. Therefore,
to experimentally test the compensation mechanism of the equation (39)
type itself, one must include configurations in which the condensate is
cut out across closure units, so that in general \[Q_C\neq0\] can occur.

However, this does not mean that

\[
Q_C=0
\]

implies the absence of any background deformation whatsoever. While the
total closure of the sum of squares is maintained, the internal
correlations, the phase arrangement, or the relational structure between
the condensate and its exterior may still be rearranged. This paper
therefore distinguishes

\begin{itemize}
\tightlist
\item
  \textbf{residual-compensation-type distortion}: the deformation
  required by equation (39) for \[Q_C\neq0\],
\item
  \textbf{closure-preserving correlational distortion}: a different kind
  of rearrangement that can occur while maintaining \[Q_C=0\].
\end{itemize}

The existence conditions and the readout of the latter are an open
problem.

\subsubsection{13.3 Spectral readout of the
distortion}\label{spectral-readout-of-the-distortion}

Inverse-transform the pre/post-condensation difference
\[\delta\mathcal A(q)\] to obtain \[\delta\rho(Y)\].

Conceptually,

\[
\delta\rho(Y)
=
\mathcal F^{-1}
\left[
\delta\mathcal A(q)
\right].
\tag{40}
\]

However, when the spatial scale is obtained from the logarithmic
amplitude, this is interpreted, with respect to the original complex
state variables, as the product of an inverse Fourier transform and an
inverse Mellin transform.

If a localized condensate truly distorts the surrounding gauge, then
\[\delta\rho(Y)\neq0\] exists outside the condensation center as well,
and

\[
\delta\rho(Y)
\to0
\qquad
\text{for sufficiently large relative separation}
\tag{41}
\]

is expected.

In that case, one obtains the structure required of an isolated object:
a condensation center, a finite core, and a background deformation
decaying outward.

\begin{center}\rule{0.5\linewidth}{0.5pt}\end{center}

\subsection{14. Relation to Background
Curvature}\label{relation-to-background-curvature}

The previous paper {[}K1{]} treated the possibility of constructing an
additive readout space from a closed, bounded complex state field.

In this paper, the state-distribution deformation caused by the
condensate is extracted as \[\delta\rho\].

Writing the readout metric as \[g_0\] and the post-condensation one as
\[g_C\], it will in the future be possible to define

\[
\boxed{
\delta g
=
g_C-g_0
}
\tag{42}
\]

as the local metric deviation caused by condensation.

However, this paper does not yet derive a unique map from \[\delta\rho\]
to \[\delta g\].

What this paper establishes is therefore up to

\[
\boxed{
\text{localized condensation}
\longrightarrow
\text{inhomogeneity of the vacuum state distribution}
}
\]

while

\[
\boxed{
\delta\rho
\longrightarrow
\text{Einstein curvature}
}
\]

is left as an unsolved problem.

This distinction is important.

On the other hand, whether a surrounding deformation due to condensation
exists is directly verifiable by numerical experiment.

\begin{center}\rule{0.5\linewidth}{0.5pt}\end{center}

\subsection{15. Relation to the Four-Dimensional
FFT}\label{relation-to-the-four-dimensional-fft}

On the surface, equation (29) is close to a four-dimensional Fourier
transform.

In the usual four-dimensional FFT, however, the lattice of
\[ (x,y,z,t) \] exists from the start.

In this paper the order is reversed:

\[
\boxed{
\text{anonymous complex waves}
\rightarrow
\text{relative scale and phase readout}
\rightarrow
\text{emergent coordinates}
\rightarrow
\text{spectral condensation analysis}
}
\tag{43}
\]

In particular, since

\[
X=\ln r,
\]

we have

\[
e^{-isX}
=
e^{-is\ln r}
=
r^{-is}.
\]

For the original amplitude variable \[r\], this is the Mellin kernel
itself.

Hence the operation of this paper is not a mere four-dimensional FFT but
carries the correspondence

\[
\boxed{
\text{Fourier in the emergent logarithmic space}
\equiv
\text{Mellin in the original amplitude space}
}
\]

\begin{center}\rule{0.5\linewidth}{0.5pt}\end{center}

\subsection{16. Basis Search from Anonymous
States}\label{basis-search-from-anonymous-states}

The equations so far give the moment extraction after readout
coordinates corresponding to the condensate have been obtained.

The hardest problem, however, is how to choose that basis itself from
the anonymous state.

This paper proposes to consider, for a candidate basis \[B\], the
spectrum

\[
\mathcal A_B(q)
\]

and to define a condensation evaluation functional

\[
\mathcal C(B).
\]

For example,

\[
\mathcal C(B)
=
\frac{
\displaystyle
\int_{\Omega_{\rm peak}}
|\delta\mathcal A_B(q)|^2dq
}{
\displaystyle
\int_{\Omega_{\rm all}}
|\delta\mathcal A_B(q)|^2dq
}
\tag{44}
\]

measures the concentration into a limited spectral region.

Define the optimal basis as

\[
\boxed{
B_*
=
\operatorname*{arg\,max}_{B}
\mathcal C(B)
}
\tag{45}
\]

However, the admissible bases \[B\] are not arbitrary: they must be
constructed only from permutation-invariant state quantities, must not
depend on the common-scale and common-phase gauges, must preserve the
square-zero closure, and must be consistent with the spacetime readout
of the previous paper {[}K1{]}.

This is therefore not an operation of freely fitting coordinates to
fabricate an arbitrary condensation.

Rather, it is the problem of

\[
\boxed{
\text{searching, within the bases the zero-closure system itself admits, for the most self-consistently localized mode}
}
\]

\begin{center}\rule{0.5\linewidth}{0.5pt}\end{center}

\subsection{17. Minimal Verification as Numerical
Experiments}\label{minimal-verification-as-numerical-experiments}

The proposal of this paper is directly verifiable numerically.

\subsubsection{Experiment A: Anonymity}\label{experiment-a-anonymity}

Create a complex wave set containing a known condensate and randomly
shuffle the numbering of all waves.

The pass condition is

\[
\mathcal A_{\rm before}
=
\mathcal A_{\rm after}.
\]

\subsubsection{Experiment B: Common-scale
invariance}\label{experiment-b-common-scale-invariance}

Transform all waves by

\[
z_n\mapsto\kappa z_n.
\]

The pass condition is

\[
|\mathcal A_{\kappa}(s,m)|^2
=
|\mathcal A(s,m)|^2.
\]

\subsubsection{Experiment C: Common-phase
invariance}\label{experiment-c-common-phase-invariance}

Transform by

\[
z_n\mapsto e^{i\phi}z_n.
\]

The pass condition is likewise

\[
|\mathcal A_{\phi}(s,m)|^2
=
|\mathcal A(s,m)|^2.
\]

\subsubsection{Experiment D: Center-of-mass
recovery}\label{experiment-d-center-of-mass-recovery}

Apply a known relative shift \[\Delta X\] to the condensate.

The prediction is

\[
\chi_C(s)
\mapsto
 e^{-is\Delta X}
\chi_C(s),
\]

and

\[
{-}
\left.
\frac{d}{ds}
\Delta\arg\chi_C
\right|_0
=
\Delta X
\tag{46}
\]

should be recovered.

\subsubsection{Experiment E: Finite-width
recovery}\label{experiment-e-finite-width-recovery}

Generate a condensate with a known variance \[\sigma^2\].

Compare the variance recovered via equation (24) with the input value.

\subsubsection{Experiment F: Multiple
condensates}\label{experiment-f-multiple-condensates}

Place two or more condensates and recover \[X_A-X_B\] from the cross
phase.

Confirm that the result is invariant under changes of the absolute
origin.

\subsubsection{Experiment G: Background
distortion}\label{experiment-g-background-distortion}

Generate a local condensation while maintaining the global zero closure,
and measure how far \[\delta\rho(Y)\] extends outside the condensation
core. In particular, for the series testing the residual-compensation
mechanism of the equation (39) type, do not select only complete closure
shells as the condensate; include in the generation conditions
condensations cut out across closure units, so that in general
\[Q_C\neq0\] can occur. In parallel, measure as a separate series
whether closure-preserving correlational distortion arises even for
condensations maintaining \[Q_C=0\].

This is the most physically important verification in this paper.

\subsubsection{Experiment H: Blind control and false-positive
rate}\label{experiment-h-blind-control-and-false-positive-rate}

Generate, as a control group, homogeneous vacuum states with no
condensate planted at all, and run the basis search of Section 16 under
the same conditions. Measure the distributions of the concentration
\[\mathcal C(B)\] for each candidate basis, the difference spectrum
\[\delta\mathcal A\], the apparent center of mass and finite width, and
the background distortion indicators, and first establish the
false-positive rate that arises even in systems containing no
condensate.

Then require that the signals of Experiments D--G be statistically
separable from this control distribution. This discriminates whether the
basis search is discovering a condensation or manufacturing an apparent
localization through optimization.

\begin{center}\rule{0.5\linewidth}{0.5pt}\end{center}

\subsection{18. Relation to
Renormalization}\label{relation-to-renormalization}

The foundational equation of this paper,

\[
\sum_n z_n^2=0,
\]

is invariant under the common scale transformation.

Therefore no particular absolute scale needs to be chosen to preserve
the self-consistency of the foundational state.

Moreover, the mass-like measure is not injected as a bare mass
parameter; it is read out as the condensation weight, the relative
center of mass, the finite width, and the background deformation.

For this reason, this construction at least does not require, as a
foundational definition, ``readjusting divergent bare parameters to
observables''.

However, we do not yet conclude from this fact alone that all
ultraviolet divergences of quantum field theory automatically disappear.

The strict claim of this paper is only that

\[
\boxed{
\text{the foundational closure condition is scale invariant, and mass-like quantities are not introduced as external bare parameters}
}
\]

If, in the future, all observables are shown to remain finite when
interactions between condensates are introduced, only then can a
stronger UV-finiteness be discussed.

\begin{center}\rule{0.5\linewidth}{0.5pt}\end{center}

\subsection{19. Difference from Standard Mass
Concepts}\label{difference-from-standard-mass-concepts}

The ``mass-like measure'' of this paper does not yet identify the Higgs
mechanism of the Standard Model, inertial mass, and gravitational mass.

What this paper constructs is

\[
\boxed{
\text{a measure of localized condensation that creates an intrinsic center of mass inside the homogeneous vacuum}
}
\]

The quantities to be distinguished and verified in the future are
therefore at least: the condensation participation weight \[M_0\], mass
as inertial response, the strength of the surrounding background
distortion, the coefficient of the gravitational far field, and the
internal oscillation period or eigenfrequency.

If these eventually converge to one and the same condensation invariant,
that would constitute a unified readout of mass.

At the present stage this coincidence is not assumed.

\begin{center}\rule{0.5\linewidth}{0.5pt}\end{center}

\subsection{20. Minimality of the
Construction}\label{minimality-of-the-construction}

The mathematical structure added in this paper is limited.

What is needed as the foundation is

\[
\boxed{
\begin{aligned}
&\text{(i) an anonymous complex state set},\\
&\text{(ii) the square-zero closure},\\
&\text{(iii) the relative logarithmic scale and phase readout of the previous paper},\\
&\text{(iv) harmonic analysis on the empirical measure}
\end{aligned}
}
\]

Particle labels, external coordinate lattices, absolute origins, and
bare masses are not introduced.

Nor is a Fourier transform taken after first specifying the center of a
localized body. The center comes out of the first moment after the
transform.

The logical order is therefore

\[
\boxed{
\text{waves}
\rightarrow
\text{condensation}
\rightarrow
\text{spectrum}
\rightarrow
\text{center of mass}
}
\]

and not the order of placing a wave packet at a known position.

\begin{center}\rule{0.5\linewidth}{0.5pt}\end{center}

\subsection{21. Open Problems}\label{open-problems}

\subsubsection{21.1 Are condensates generated
spontaneously?}\label{are-condensates-generated-spontaneously}

This paper gives a method to read out condensates, but it does not prove
that

\[
\sum_n z_n^2=0
\]

alone necessarily generates stable localized condensation.

This is the most important dynamical problem.

\subsubsection{21.2 Stability of
condensates}\label{stability-of-condensates}

The conditions under which the center of mass \[\mu_C\] and the
covariance \[\Sigma_C\] remain stable over long times must be found.

\subsubsection{21.3 Relation between condensation weight and background
distortion}\label{relation-between-condensation-weight-and-background-distortion}

Whether a simple proportionality law exists between

\[
M_0
\]

and the far-field amplitude of

\[
\delta\rho
\]

must be verified.

\subsubsection{21.4 Inverse-square law}\label{inverse-square-law}

After the three-dimensional readout, whether the background distortion
of an isolated condensate produces a gradient of the

\[
\frac{1}{R^2}
\]

type at large distance is unverified.

\subsubsection{21.5 Coincidence of inertia and
gravity}\label{coincidence-of-inertia-and-gravity}

If the condensate's resistance to acceleration and the surrounding
background distortion are determined by the same invariant, this could
connect to the coincidence of inertial and gravitational mass. At the
present stage this is a prediction problem.

\subsubsection{21.6 Uniqueness of condensate
decomposition}\label{uniqueness-of-condensate-decomposition}

This paper separates the condensate component \[\mu_C\] from the
difference measure \[\delta\mu\] and reads its center of mass and
extent. However, when multiple localized correlations coexist in the
anonymous many-body system, it is not yet shown that the decomposition

\[
\delta\mu
\longrightarrow
\{\mu_{C_1},\mu_{C_2},\ldots\}
\]

is uniquely determined. Beyond the maximal concentration in the basis
search of Section 16, the existence and uniqueness of the decomposition,
the separation limit for nearby condensates, and the resolution
dependence must be formalized. This is the central problem for not
specifying from outside ``how much counts as one object''.

\subsubsection{21.7 Residual-compensation-type versus closure-preserving
background
distortion}\label{residual-compensation-type-versus-closure-preserving-background-distortion}

Under which conditions the two kinds of background deformation
distinguished in Section 13.2 arise, and whether they converge to the
same far-field readout, is unresolved. In particular, it must be
discriminated whether long-range distortion can arise from correlational
rearrangement alone even for condensates with \[Q_C=0\].

\begin{center}\rule{0.5\linewidth}{0.5pt}\end{center}

\subsection{22. Discussion --- No Separate Entity at the Vacuum-Matter
Boundary}\label{discussion-no-separate-entity-at-the-vacuum-matter-boundary}

The characteristic of this construction is that vacuum and matter are
not placed as fundamentally different things.

The homogeneous vacuum is

\[
\mathcal V_0,
\]

and the universe containing condensates is

\[
\mathcal V.
\]

Both consist of the same complex state degrees of freedom; the
difference lies in the correlation structure.

Therefore the reading

\[
\boxed{
\text{matter}
=
\text{localized condensation of the vacuum state gauge field}
}
\tag{47}
\]

becomes possible.

Furthermore, if the compensating deformation required to maintain the
global zero closure extends to the surroundings upon condensation, then
matter and background distortion are not independent inputs but two
readouts of the same self-consistent rearrangement.

Here lies the most important physical possibility of this construction.

\begin{center}\rule{0.5\linewidth}{0.5pt}\end{center}

\subsection{23. Conclusion}\label{conclusion}

In this paper, for the case where a localized condensate exists inside
the homogeneous vacuum state field of anonymous complex standing waves,

\[
\sum_n z_n^2=0,
\]

we constructed a framework that reads out the condensate's material
information without external coordinates or known masses.

The main results are as follows.

First, for the empirical measure of the anonymous complex set we defined
the Fourier--Mellin-type transform

\[
\mathcal A(s,m)
=
\frac1N
\sum_n
\exp
\left[
-i
\left(
s\ln r_n+m\theta_n
\right)
\right].
\]

Second, its spectral intensity is invariant under wave permutations, the
common amplitude scale, and the common phase rotation.

Third, from the logarithmic expansion of the condensate's normalized
transform,

\[
\ln\chi_C(q)
=
-iq^{\mathsf T}\mu_C
-
\frac12q^{\mathsf T}\Sigma_Cq
+
\cdots,
\]

we obtained the relative center of mass as the first-order term and the
finite extent and shape as the second-order term.

Fourth, by taking the unnormalized zeroth-order term as the condensation
participation weight, we constructed the mass-like moment hierarchy

\[
\boxed{
0\text{th}
\to
\text{condensation weight},
\qquad
1\text{st}
\to
\text{center of mass},
\qquad
2\text{nd}
\to
\text{extent}
}
\]

Fifth, from the pre/post-condensation difference spectrum

\[
\delta\mathcal A
=
\mathcal A-\mathcal A_0,
\]

we showed a method to directly detect the distortion of the vacuum state
gauge extending outside the condensate.

In this construction, therefore,

\[
\boxed{
\text{homogeneous vacuum}
\rightarrow
\text{localized condensation}
\rightarrow
\text{relative center of mass}
\rightarrow
\text{finite extent}
\rightarrow
\text{mass-like measure}
\rightarrow
\text{surrounding gauge distortion}
}
\]

can be described continuously as internal readouts of one and the same
anonymous complex state field.

The most important open problem is not the readout method but whether
the dynamics of a square-zero-closed homogeneous standing-wave vacuum
can be shown to spontaneously generate such stable condensates of finite
width.

In that sense this paper is not a ``completed theory of mass''; it
provides

\[
\boxed{
\text{the minimal mathematical apparatus for reading an object as an object out of the anonymous sea of vacuum waves}
}
\]

\begin{center}\rule{0.5\linewidth}{0.5pt}\end{center}

\section{Appendix A: Minimal Extension to the Self-Dual Vacuum and
Non-Null
Readouts}\label{appendix-a-minimal-extension-to-the-self-dual-vacuum-and-non-null-readouts}

\subsection{A.1 Purpose and
positioning}\label{a.1-purpose-and-positioning}

In the countably infinite special solution of the previous paper
{[}K1{]}, the spatial scale and the period scale are proportional to the
same shell parameter, and the post-mapping wavelength-frequency product
is constant for every shell pair. Moreover, when the readout units are
matched, all shell-pair intervals become null.

In this appendix, we generalize this structure minimally without
breaking it, and investigate under which conditions the three-way
classification into spacelike, null, and timelike opens up. What is
derived here is a mathematical classification between readout scales;
that localized condensates actually generate timelike structure is
distinguished at the end as an independent hypothesis.

Let the positive spatial and period scales be

\[
L_k>0,
\qquad
P_k>0,
\]

with logarithmic readouts

\[
X_k=C_X\ln L_k,
\qquad
T_k=C_T\ln P_k.
\tag{A1}
\]

For a shell pair \[j,k\],

\[
\Delta X_{jk}
=
C_X\ln\frac{L_j}{L_k},
\qquad
\Delta T_{jk}
=
C_T\ln\frac{P_j}{P_k}.
\tag{A2}
\]

Define the wavelength readout and the period/frequency readouts as

\[
\lambda_{jk}=|\Delta X_{jk}|,
\qquad
\tau_{jk}=|\Delta T_{jk}|,
\qquad
\nu_{jk}=\frac{1}{\tau_{jk}}.
\tag{A3}
\]

\begin{center}\rule{0.5\linewidth}{0.5pt}\end{center}

\subsection{\texorpdfstring{A.2 The necessary and sufficient condition
for constant \(\lambda\nu\) across all shell
pairs}{A.2 The necessary and sufficient condition for constant \textbackslash lambda\textbackslash nu across all shell pairs}}\label{a.2-the-necessary-and-sufficient-condition-for-constant-lambdanu-across-all-shell-pairs}

\subsubsection{Theorem A1}\label{theorem-a1}

Consider a shell family containing at least three mutually distinct
spatial scales \[L_k\]. For all distinct shell pairs \[j\neq k\],

\[
\lambda_{jk}\nu_{jk}=K
\qquad
(K>0)
\tag{A4}
\]

is one and the same constant if and only if there exist a constant
\[P_0>0\] and a nonzero real number \[\alpha\] such that

\[
\boxed{
P_k=P_0L_k^{\alpha}
}
\tag{A5}
\]

holds for all shells. Since a reversal of axis orientation can be
absorbed into the sign of \[\alpha\] or the sign of \[C_T\], one may
take \[\alpha>0\] when treating only increasing scale relations.

\subsubsection{Proof}\label{proof}

Set

\[
x_k=\ln L_k,
\qquad
y_k=\ln P_k.
\]

From equations (A3)(A4),

\[
\frac{|C_X|\,|x_j-x_k|}
{|C_T|\,|y_j-y_k|}
=K,
\]

so for some \[a>0\],

\[
|y_j-y_k|
=
a|x_j-x_k|
\tag{A6}
\]

holds for all shell pairs.

Therefore the map \[x_k\mapsto y_k/a\] preserves all pairwise distances
of this point set on the real line. Since a distance-preserving map on a
subset of the real line containing at least three points is unique up to
translation and orientation reversal, there exist \[b\in\mathbb R\] and
\[\sigma=\pm1\] such that

\[
y_k
=
\sigma a x_k+b.
\tag{A7}
\]

Setting

\[
\alpha=\sigma a,
\qquad
P_0=e^b
\]

yields equation (A5).

Conversely, if equation (A5) holds, then

\[
\ln\frac{P_j}{P_k}
=
\alpha
\ln\frac{L_j}{L_k},
\]

so equations (A2)(A3) give

\[
\boxed{
\lambda_{jk}\nu_{jk}
=
\frac{|C_X|}{|C_T|\,|\alpha|}
}
\tag{A8}
\]

which does not depend on the shell pair. This establishes necessity and
sufficiency. \(\square\)

By this theorem, the condition \[P_k\propto L_k\] of the previous paper
{[}K1{]} is the special point

\[
\boxed{\alpha=1}
\]

of the general condition.

\begin{center}\rule{0.5\linewidth}{0.5pt}\end{center}

\subsection{\texorpdfstring{A.3 Spacelike / null / timelike
classification by
\(\alpha\)}{A.3 Spacelike / null / timelike classification by \textbackslash alpha}}\label{a.3-spacelike-null-timelike-classification-by-alpha}

Normalize the readout units to

\[
|C_X|=|C_T|=:C.
\tag{A9}
\]

From equation (A5),

\[
\Delta T_{jk}
=
\sigma\alpha\,\Delta X_{jk},
\]

where \[\sigma=\pm1\] expresses only the axis orientation and does not
affect the quadratic form.

Taking the Lorentz-type readout quadratic form

\[
ds_{jk}^2
=
(\Delta X_{jk})^2
-
(\Delta T_{jk})^2,
\tag{A10}
\]

we obtain

\[
\boxed{
 ds_{jk}^2
=
(1-\alpha^2)
C^2
\left(
\ln\frac{L_j}{L_k}
\right)^2
}
\tag{A11}
\]

Therefore, for \[\alpha>0\],

\[
\boxed{
0<\alpha<1
\quad\Longrightarrow\quad
 ds^2>0
\quad\text{(spacelike)}
}
\tag{A12}
\]

\[
\boxed{
\alpha=1
\quad\Longrightarrow\quad
 ds^2=0
\quad\text{(null)}
}
\tag{A13}
\]

\[
\boxed{
\alpha>1
\quad\Longrightarrow\quad
 ds^2<0
\quad\text{(timelike)}
}
\tag{A14}
\]

Hence the special solution of the previous paper is the self-dual point

\[
\boxed{
\alpha=1
}
\]

at which the exponents of the spatial and period scales coincide, and
there all shell-pair intervals are null.

Also, from equation (A8), in the same normalization,

\[
\boxed{
\lambda_{jk}\nu_{jk}
=
\frac{1}{\alpha}
}
\tag{A15}
\]

In a normalization carrying a dimensional reference speed \[c_*\], the
right-hand side becomes \[c_*/\alpha\].

What matters is that the constancy of \[\lambda\nu\] itself is not
restricted to null. Constancy across all shell pairs is equivalent to
the power relation \[P\propto L^\alpha\], and within it \[\alpha=1\]
selects the null self-dual point.

\begin{center}\rule{0.5\linewidth}{0.5pt}\end{center}

\subsection{A.4 A rational-exponent shell family preserving integrality
and exact
closure}\label{a.4-a-rational-exponent-shell-family-preserving-integrality-and-exact-closure}

We now construct the \[\alpha\] generalization above as a countable
complex zero-closure solution with no rounding error.

Take the base integer sequence

\[
B_k=k(k+1),
\qquad
k=1,2,3,\ldots
\tag{A16}
\]

and for positive integers \[p,q\] define

\[
R_k=B_k^q,
\qquad
N_k=B_k^p.
\tag{A17}
\]

Take the complex waves of each shell as

\[
\boxed{
z_{k,n}
=
\frac{A}{R_k}
\exp\left[
i\left(
\phi_k+
\frac{\pi n}{N_k}
\right)
\right],
\qquad
n=0,\ldots,N_k-1
}
\tag{A18}
\]

Squaring,

\[
z_{k,n}^2
=
\frac{A^2e^{2i\phi_k}}{R_k^2}
\exp\left(
\frac{2\pi i n}{N_k}
\right),
\]

so for \[N_k\ge2\] the sum of roots of unity makes each shell exactly
satisfy

\[
\boxed{
\sum_{n=0}^{N_k-1}z_{k,n}^2=0
}
\tag{A19}
\]

Meanwhile, \[R_k\] can be read from the inverse-amplitude scale, and
\[N_k\] from the phase-cycle order or the shell multiplicity. From
equation (A17),

\[
\ln N_k
=
\frac{p}{q}
\ln R_k
\tag{A20}
\]

holds exactly.

Therefore we obtain the rational-exponent family

\[
\boxed{
\alpha=\frac{p}{q}
}
\tag{A21}
\]

The previous paper's \[\alpha=1\] corresponds to \[p=q\].

\subsubsection{Absolute convergence
condition}\label{absolute-convergence-condition}

The total sum of squared absolute values over all shells is

\[
\sum_k\sum_{n=0}^{N_k-1}|z_{k,n}^2|
=
|A|^2
\sum_k
\frac{N_k}{R_k^2},
\]

and substituting equation (A17),

\[
|A|^2
\sum_k
B_k^{p-2q}.
\tag{A22}
\]

Since \[B_k\sim k^2\], the convergence condition is

\[
2(p-2q)<-1.
\]

As \[p,q\] are integers, this is equivalent to

\[
\boxed{
p\le2q-1
}
\tag{A23}
\]

Therefore, in the absolutely convergent rational-exponent family,

\[
\boxed{
0<\alpha
=
\frac pq
\le
2-
\frac1q
<2
}
\tag{A24}
\]

This range contains all of

\begin{itemize}
\tightlist
\item
  \[p<q\]: spacelike,
\item
  \[p=q\]: null,
\item
  \[q<p\le2q-1\]: timelike.
\end{itemize}

In particular, having an exact timelike family requires \[q\ge2\].

Hence, without breaking the null solution of the previous paper, the
construction extends to both spacelike and timelike while simultaneously
maintaining integer cycle numbers, exact per-shell square-zero closure,
and overall absolute convergence.

\begin{center}\rule{0.5\linewidth}{0.5pt}\end{center}

\subsection{A.5 Connection to condensates --- separating derived results
from
hypotheses}\label{a.5-connection-to-condensates-separating-derived-results-from-hypotheses}

In the main text, the localized condensate was defined as a local
rearrangement of the homogeneous vacuum state field, and the
condensation weight, the relative center of mass, the finite extent, and
the background state gauge distortion were constructed from the
Fourier--Mellin-type readout.

What this appendix newly derives is the following.

\begin{enumerate}
\def\labelenumi{\arabic{enumi}.}
\tightlist
\item
  The constancy of \[\lambda\nu\] across all shell pairs is necessarily
  and sufficiently equivalent to \[P\propto L^\alpha\].
\item
  \[\alpha=1\] is the self-dual point at which all shell pairs are null.
\item
  \[\alpha<1\] and \[\alpha>1\] give spacelike and timelike readouts,
  respectively.
\item
  For rational exponents \[\alpha=p/q\], there exists a countable shell
  family preserving exact square-zero closure and absolute convergence.
\end{enumerate}

Furthermore, since exchanging the roles of \[L\] and \[P\] formally
gives \[\alpha\leftrightarrow1/\alpha\], timelike-type and
spacelike-type readouts may form a reciprocal pair across the self-dual
point \[\alpha=1\]. This is reminiscent of the reciprocal structure of
phase velocity and group velocity appearing for standard massive waves,
but this paper does not derive the physical identity of the two. It
therefore remains a \textbf{correspondence hypothesis}.

When testing this hypothesis, one must distinguish whether
\[\alpha_{\mathrm{eff}}\] is measured from ``the worldline-type readout
of the condensate'' or from ``the wavefront-type readout attached to the
condensate''. The two are not to be unconditionally identified as the
same \[\alpha_{\mathrm{eff}}\].

On the other hand, the following proposition is not yet derived in this
paper.

\begin{quote}
\textbf{Connection hypothesis}: the formation of a localized condensate
may locally break the \[\alpha=1\] self-dual relation of the homogeneous
vacuum and generate a timelike readout with \[\alpha_{\mathrm{eff}}>1\]
inside or around the condensate.
\end{quote}

The local effective exponent can be formally defined, for two nearby
readout scales, as

\[
\boxed{
\alpha_{\mathrm{eff}}
=
\frac{\Delta\ln P}{\Delta\ln L}
}
\tag{A25}
\]

If

\[
\alpha_{\mathrm{eff}}=1
\]

is maintained in the homogeneous vacuum and systematically shifts to

\[
\alpha_{\mathrm{eff}}>1
\]

near the condensation center, a direct relation between mass-like
condensation and timelike readout would be suggested.

At the present stage, however,

\[
\boxed{
\text{localized condensation}
\Longrightarrow
\alpha_{\mathrm{eff}}>1
}
\]

is a hypothesis and is not treated as a consequence.

\begin{center}\rule{0.5\linewidth}{0.5pt}\end{center}

\subsection{A.6 Minimal discriminating
computations}\label{a.6-minimal-discriminating-computations}

This connection hypothesis can be discriminated in the following order.

\subsubsection{A.6.1 Direct verification of the rational-exponent
family}\label{a.6.1-direct-verification-of-the-rational-exponent-family}

Generate equation (A18) for several \[p,q\] and confirm that

\begin{itemize}
\tightlist
\item
  each shell's sum of squares is zero within numerical precision,
\item
  \[\ln N_k/\ln R_k=p/q\],
\item
  \[\lambda\nu\] is constant across all shell pairs,
\item
  the sign of equation (A11) matches the order relation of \[p/q\].
\end{itemize}

\subsubsection{A.6.2 Readout stability against local closure
residuals}\label{a.6.2-readout-stability-against-local-closure-residuals}

Give two shells \[j,k\] several small phase degrees of freedom and
construct exactly

\[
Q_j=\epsilon,
\qquad
Q_k=-\epsilon,
\qquad
Q_{\mathrm{total}}=0.
\tag{A26}
\]

For a phase deformation of a single element,

\[
\delta z
\simeq
iz\,\delta\theta,
\]

we have

\[
\delta(z^2)
\simeq
2iz^2\,\delta\theta,
\tag{A27}
\]

so the complex direction of the residual is constrained to \[z^2\], and
a single degree of freedom is in general insufficient to produce an
arbitrary complex \[\epsilon\]. Therefore prepare at least two
independent real phase degrees of freedom in each shell and solve the
two conditions for the real and imaginary parts.

On top of that, measure whether the Fourier--Mellin readouts of

\begin{itemize}
\tightlist
\item
  the condensation peak position,
\item
  the relative center of mass,
\item
  the finite extent,
\item
  \[\alpha_{\mathrm{eff}}\],
\end{itemize}

remain continuously stable as \[\epsilon\to0\]. For
\[\alpha_{\mathrm{eff}}\], always record whether the value is computed
from the worldline-type readout or the wavefront-type readout, and
independently discriminate whether the two stand in a reciprocal
relation.

If this discriminating computation is stable, then even though the exact
special solutions are of measure zero, the likelihood increases that
physically readable structure exists in their neighborhood.

\begin{center}\rule{0.5\linewidth}{0.5pt}\end{center}

\subsection{A.7 Conclusion of the
appendix}\label{a.7-conclusion-of-the-appendix}

The \[\alpha=1\] special solution of the previous paper {[}K1{]} is not
merely one example where \[\lambda\nu\] is constant; it is the self-dual
point of the power relation

\[
P\propto L^\alpha,
\]

at which all shell pairs become null.

In this appendix we showed the necessary and sufficient condition

\[
\boxed{
\lambda\nu=\mathrm{const}
\quad\Longleftrightarrow\quad
P=P_0L^\alpha
}
\]

and, via the exact rational-exponent shell family, constructed

\[
\boxed{
\alpha<1:
\text{spacelike},
\qquad
\alpha=1:
\text{null},
\qquad
\alpha>1:
\text{timelike}
}
\]

within the same square-zero closure system.

Therefore the null structure is not forced by the square-zero closure
itself; it is selected by the self-duality condition \[\alpha=1\]
between the spatial scale and the period scale.

Whether localized condensation produces a local deviation from this
self-dual point, and whether that deviation connects to the mass-like
measure and the timelike readout, is the next discrimination problem
obtained directly from this paper.

\begin{center}\rule{0.5\linewidth}{0.5pt}\end{center}

\subsection{References}\label{references}

\subsubsection{Self-citation}\label{self-citation}

{[}K1{]} Noriaki Kihara, \textbf{``Emergent Spacetime and Wave-Scale
Readout from Anonymous Complex Zero Closure --- Minimal Readout, an
Infinite Constructive Solution, Logarithmic Spacetime Mapping, and
Post-Mapping \(\lambda\nu\) Invariance''}, 2026, v2.0. Zenodo Concept
DOI: 10.5281/zenodo.22282217.

\subsubsection{External references}\label{external-references}

{[}E1{]} D. Casasent and D. Psaltis, ``Position, rotation, and scale
invariant optical correlation,'' \emph{Applied Optics}, \textbf{15}(7),
1795--1799 (1976). DOI: 10.1364/AO.15.001795.

{[}E2{]} Eugene Lukacs, \emph{Characteristic Functions}, 2nd ed., Hafner
Publishing Company / Griffin, 1970.

\begin{center}\rule{0.5\linewidth}{0.5pt}\end{center}

\subsection{Note: The Original Scope of This
Paper}\label{note-the-original-scope-of-this-paper}

Reference {[}E1{]} is cited for the known status of Fourier--Mellin-type
scale- and rotation-invariant analysis, and {[}E2{]} for the standard
mathematics of obtaining moments from characteristic functions and
cumulant expansions.

These references do not propose the anonymous complex vacuum state field
satisfying

\[
\sum_n z_n^2=0
\]

of this paper, the definition of localized condensates within it, the
construction reading the condensation center of mass as a mass-like
measure, or the vacuum-state gauge distortion caused by condensation.

The original proposal of this paper is to connect these known
mathematics to the emergent spacetime readout from anonymous complex
zero closure of the previous paper {[}K1{]}, and to construct

\[
\boxed{
\text{anonymous complex states}
\rightarrow
\text{condensate}
\rightarrow
\text{relative center of mass}
\rightarrow
\text{mass-like measure}
\rightarrow
\text{background gauge distortion}
}
\]

as one readout sequence.

\end{document}
